\notechapter{N-20220320}{A Review of Calculus: Mean Value Theorems and Continuity}

\section{From the mean value theorem to Cauchy's form}

This note continues the calculus review.  The main topic is the mean value theorem and its Cauchy form.

\begin{theorem}[Mean value theorem]
Let \(f:[a,b]\to\mathbb R\) be continuous on \([a,b]\) and differentiable on \((a,b)\).  Then there exists \(c\in(a,b)\) such that
\[
  f'(c)=\frac{f(b)-f(a)}{b-a}.
\]
\end{theorem}

Geometrically, there is a point at which the tangent slope equals the slope of the secant line through the endpoints.

\begin{theorem}[Cauchy's mean value theorem]
Let \(f,g:[a,b]\to\mathbb R\) be continuous on \([a,b]\) and differentiable on \((a,b)\).  If \(g'(x)\neq0\) on \((a,b)\), then there exists \(c\in(a,b)\) such that
\[
  \frac{f'(c)}{g'(c)}=\frac{f(b)-f(a)}{g(b)-g(a)}.
\]
Equivalently,
\[
  (f(b)-f(a))g'(c)=(g(b)-g(a))f'(c).
\]
\end{theorem}

\begin{remark}
This is the version sometimes called Cauchy's mean value theorem.  It is the natural two-function generalization of the ordinary mean value theorem and is often the cleanest route to l'Hospital-type arguments.
\end{remark}

\section{Continuity at a point}

The note then returns to continuity.

\begin{definition}
Let \(f\) be a real-valued function defined in a neighborhood of \(a\in\mathbb R\).  The function \(f\) is continuous at \(a\) if
\[
  \lim_{x\to a}f(x)=f(a).
\]
Equivalently, for every \(\varepsilon>0\) there is \(\delta>0\) such that
\[
  |x-a|<\delta\Rightarrow |f(x)-f(a)|<\varepsilon.
\]
\end{definition}

\begin{remark}
The intuitive statement is that when \(x\) is close to \(a\), the value \(f(x)\) is close to \(f(a)\).  The epsilon-delta form is the precise version of that intuition.
\end{remark}

\section{How this note fits the previous one}

The previous review note introduced limits, derivatives, big-O/small-o notation, and integration.  This continuation adds the mean value theorem and continuity, which are the bridge between local derivative information and global information about a function on an interval.

The original text layer for this note was short and duplicated in places, so the missing proof details are filled in explicitly below.  The key point is that both the ordinary mean value theorem and Cauchy's mean value theorem are consequences of Rolle's theorem: one subtracts a secant line in the first case, and a suitable linear combination of two functions in the second.


\section{Rolle's theorem as the engine}

The mean value theorem is usually proved from Rolle's theorem.  Rolle says that if \(f(a)=f(b)\), and \(f\) is continuous on \([a,b]\) and differentiable on \((a,b)\), then there is \(c\in(a,b)\) such that
\[
  f'(c)=0.
\]
Geometrically, a curve beginning and ending at the same height must have a horizontal tangent somewhere.

To prove the mean value theorem, subtract the secant line.  Define
\[
  h(x)=f(x)-\frac{f(b)-f(a)}{b-a}(x-a).
\]
Then \(h(a)=h(b)\).  Rolle's theorem gives \(h'(c)=0\), which is exactly
\[
  f'(c)=\frac{f(b)-f(a)}{b-a}.
\]
This proof explains the theorem better than the formula alone: the secant line is removed so that the endpoints line up.

\section{Cauchy's theorem and L'Hopital's rule}

Cauchy's mean value theorem applies Rolle's theorem to a linear combination of two functions.  It is the natural source of L'Hopital's rule.  Suppose
\[
  \lim_{x\to a}f(x)=\lim_{x\to a}g(x)=0
\]
and \(g'(x)\ne0\).  Applying Cauchy's theorem to \(f\) and \(g\) on \([a,x]\) gives a point \(c\) between \(a\) and \(x\) such that
\[
  \frac{f(x)-f(a)}{g(x)-g(a)}=\frac{f'(c)}{g'(c)}.
\]
If \(f'(c)/g'(c)\) tends to a limit as \(x\to a\), then so does \(f(x)/g(x)\).  This derivation shows why L'Hopital's rule has hypotheses; it is not a symbolic cancellation rule.

\section{Uniform continuity}

Continuity at each point allows \(\delta\) to depend on both \(\varepsilon\) and the point.  Uniform continuity allows \(\delta\) to depend only on \(\varepsilon\).  On a compact interval, every continuous function is uniformly continuous.

This fact is often proved by contradiction using sequences, or by compactness using finite subcovers.  It is the hidden reason many limiting arguments on closed intervals work smoothly.

\section{Expanded synthesis and advanced context}

Rolle, the mean value theorem, and Cauchy's mean value theorem are not isolated formulas.  They are mechanisms for turning global information on an interval into local derivative information at some point.  This is the central bridge of one-variable calculus: endpoint data forces a tangent slope somewhere in between.



\paragraph{Expanded synthesis.}
The common theme of this chapter is that an elementary limiting or computational rule becomes powerful only after one specifies the space in which the limiting process takes place. The earlier sections (`From the mean value theorem to Cauchy's form`, `Continuity at a point`, `How this note fits the previous one`, `Rolle's theorem as the engine`, `Cauchy's theorem and L'Hopital's rule`) may look like separate techniques, but they all ask the same question: which operations are stable under approximation? A derivative is stable first-order approximation,
\[
  f(a+h)=f(a)+L(h)+o(|h|),
\]
while an integral is stable accumulation with respect to a measure, and a convergent series is a limit in a chosen norm or topology.

The advanced bridge is functional-analytic. Uniform convergence is convergence in a sup norm; $L^p$ convergence is convergence in an integral norm; differentiability is the existence of a continuous linear approximation; many differential equations are operator equations $L[u]=f$. Theorems such as dominated convergence, the inverse function theorem, the projection theorem in Hilbert spaces, and semigroup existence results are refined answers to the same elementary concern: when may we pass from finite approximations to a genuine limiting object?

To use this viewpoint when rereading the original examples, identify the hidden stability condition. A Taylor expansion asks for control of the remainder. A change of variables asks for differentiability and Jacobian control. Termwise differentiation of a series asks for uniform convergence of derivative series. A differential-equation formula asks whether the operator is invertible under the imposed initial or boundary data. The elementary computation is the visible part; the advanced theory names the hypotheses that make it legitimate.


\section{Elementary-to-advanced bridge: mean value theorem as compactness plus connectedness}

Rolle's theorem and the mean value theorem are often memorized as calculus facts, but their proof rests on two deeper one-dimensional principles.  First, a continuous function on a closed interval attains a maximum and a minimum; this is compactness.  Second, differentiability turns information at an extremum into the equation $f'(\xi)=0$.

For the usual mean value theorem, define
\[
  g(x)=f(x)-\frac{f(b)-f(a)}{b-a}(x-a).
\]
Then $g(a)=g(b)$, so Rolle's theorem gives $g'(\xi)=0$, hence
\[
  f(b)-f(a)=f'(\xi)(b-a).
\]
The theorem is therefore a statement that a function and its secant line must have a parallel tangent somewhere.

\section{Elementary-to-advanced bridge: Lipschitz estimates and uniqueness}

A powerful corollary is the Lipschitz estimate.  If $|f'(x)|\le M$ on an interval, then
\[
  |f(x)-f(y)|\le M|x-y|.
\]
This turns derivative bounds into metric bounds.

This idea reappears in differential equations.  Suppose $y'=F(t,y)$ and $F$ is Lipschitz in $y$:
\[
  |F(t,u)-F(t,v)|\le L|u-v|.
\]
If $y_1,y_2$ are two solutions with the same initial value, then
\[
  |y_1(t)-y_2(t)|\le \int_{t_0}^t L|y_1(s)-y_2(s)|\,ds.
\]
Gronwall's inequality forces $y_1=y_2$.  Thus a tangent-slope theorem becomes the analytic mechanism behind uniqueness of solutions.

\section{Elementary-to-advanced bridge: Taylor formula with remainder}

The mean value theorem also generates Taylor's theorem.  For example,
\[
  f(x)=f(a)+f'(a)(x-a)+\frac{f''(\xi)}2(x-a)^2
\]
for some $\xi$ between $a$ and $x$.  The remainder measures the error of replacing a function by its local polynomial model.  In numerical analysis this error term controls approximation; in geometry it is the beginning of the language of jets.
